\documentclass[11pt]{article}
\usepackage{amsmath}
\usepackage{booktabs}
\usepackage{geometry}
\geometry{margin=1in}

\title{An Explainable Multi-Agent Framework for Budget-Constrained Travel Itinerary Optimization}
\author{TravelAgenticAI Local Demo}
\date{Draft generated from implemented repository artifacts}

\begin{document}
\maketitle

\begin{abstract}
This draft describes a local-first multi-agent travel planning system that combines typed constraint extraction,
transparent provider abstractions, deterministic budget accounting, geospatial scoring, weather-aware scheduling,
and critic validation. The implementation avoids paid APIs by default and labels estimated, fallback, synthetic,
open-data-inspired, and live data sources separately. Experimental results must be regenerated from the repository
scripts before making empirical claims.
\end{abstract}

\textbf{Keywords:} multi-agent systems, itinerary optimization, explainable AI, budget constraints, local-first software.

\section{Introduction}
Travel itinerary planning requires balancing preferences, cost, time, geography, weather, and feasibility. This project
implements an explainable decision-support system rather than a booking engine.

\section{Problem Statement}
Given a trip request $r$, candidate activities $A$, transport options $T$, accommodation options $H$, and weather
guidance $W$, select a schedule $S$ that satisfies hard constraints and maximizes a weighted utility objective.

\section{Related Work}
Verified scholarly citations should be added here before publication. This draft intentionally does not fabricate
citations.

\section{System Architecture}
The system uses a FastAPI backend, SQLite persistence, a typed orchestration graph, and a Vite React frontend. The
browser receives backend-supplied coordinates and never geocodes itinerary text.

\section{Agent Design}
The implemented flow includes intent validation, destination research, transport estimation, accommodation estimation,
weather alignment, optimization, budget reconciliation, writer, critic, and revision components.

\section{Optimization Formulation}
The deterministic objective is:
\[
\max_S \lambda_p P(S) + \lambda_b B(S) + \lambda_d D(S) + \lambda_w W(S) + \lambda_v V(S) + \lambda_h H(S)
\]
subject to:
\[
Cost(S) \leq Budget(r),\quad ActivitiesPerDay(S_d) \leq Cap(r),\quad DuplicateActivities(S)=0.
\]
The current implementation uses a transparent heuristic optimizer with deterministic budget pruning. OR-Tools CP-SAT
is a documented future improvement.

\section{Data Sources}
The default demo uses curated local datasets and deterministic estimates. Optional Open-Meteo live weather can be
enabled through environment configuration.

\section{Experimental Methodology}
Run \texttt{research/experiments/run\_benchmarks.py} followed by \texttt{run\_ablations.py}. The benchmark cases are
synthetic and transparent; no user satisfaction claims are made.

\section{Results}
Populate this section from \texttt{research/results/benchmark\_results.csv} after running the scripts.

\section{Ablation Study}
Populate this section from \texttt{research/results/ablation\_summary.csv}.

\section{Discussion}
The system emphasizes data honesty, explainability, and local reproducibility over live booking completeness.

\section{Threats to Validity}
The benchmark set is synthetic, providers are simplified, and estimated prices are not live market availability.

\section{Ethical and Privacy Considerations}
The default app stores trip requests locally in SQLite and does not require paid API credentials.

\section{Limitations}
No booking, payment, real-time inventory, or verified user study is implemented.

\section{Future Work}
Future work includes richer public-data connectors, OSRM routing, CP-SAT optimization, and verified scholarly comparison.

\section{Conclusion}
The repository demonstrates a complete local-first architecture for explainable travel itinerary decision support.

\bibliographystyle{plain}
\bibliography{references}
\end{document}
